% Exemple d'ordre du jour
% par Pierre Tremblay, Universite Laval
% modifie par Christian Gagne, Universite Laval
% 2011/01/14 - version 1.4
% modifié par Robert Bergevin, Université Laval
% 24/11/2011
% modifié par Jean-Yves Chouinard, Université Laval
% 2016/01/11
% modifié par Jean-Yves Chouinard, Université Laval
% 2017/01/04
%

%--------------------------------------------------------------------------------------
%------------------------------------- preambule --------------------------------------
%--------------------------------------------------------------------------------------
\documentclass[12pt]{ULojpv}
% Chargement des packages supplementaires
\usepackage[utf8]{inputenc}
\usepackage[T1]{fontenc}
\usepackage[french]{babel}

% Definitions des parametres de l'en-tete
\Cours{GEL--1001 Design I (méthodologie)}             % Nom du cours
\NumeroEquipe{2}                                     % Numero de l'equipe
\NomEquipe{Nord Faune}                               % Nom de l'equipe
\Objet{Ordre du jour}                                 % Nom du document
\SujetRencontre{Organisation pour le rapport version 2}        % Sujet de la rencontre
\DateRencontre{2026/02/26}                            % Date de la rencontre
\LocalRencontre{En ligne (Via Teams)}                            % Local de la rencontre
\HeureRencontre{9h30-11h20}                          % Heure de la rencontre

%--------------------------------------------------------------------------------------
%--------------------------------- corps du document ----------------------------------
%--------------------------------------------------------------------------------------
\begin{document}
\entete
\begin{enumerate}
   \item \textbf{Ouverture de la réunion}
   \item \textbf{Nomination ou confirmation du président et du secrétaire}
   \item \textbf{Lecture et adoption de l'ordre du jour}
   \item \textbf{Lecture et adoption du procès-verbal de la réunion du 19 février 2026}
   \item \textbf{Affaires découlant du procès-verbal}
      \begin{enumerate}
         \item Relecture du tableau du cahier des charges du "Rapport projet - version 1" afin d'organiser la rédaction du rapport version 2
      \end{enumerate}
   \item \textbf{Points à traiter}
      \begin{enumerate}
         \item Retour sur le séminaire du 25 février "Méthodologie du design - Conceptualisation et analyse de faisabilité  (première partie)"
         \item Organisation de la remise "Documents de Gestion G09" du 11 mars 2026
          \begin{enumerate}
               \item Répartition des tâches (Ordre du jour, procès verbal et diagramme de Gantt)
               \item Planification des livrables
            \end{enumerate}
        \item Organisation de la remise du "Rapport de Projet - version 2"  
        \begin{enumerate}
               \item Répartition des tâches (Rédaction, Révision, Remise, Correction,...)
               \item Planification des livrables
            \end{enumerate}
        
      \end{enumerate}
   \item \textbf{Divers}
   \item \textbf{Répartition des tâches}
   \item \textbf{Évaluation de la réunion}
   \item \textbf{Date, heure, lieu et objectif de la prochaine réunion}
   \item \textbf{Fermeture de la réunion}
\end{enumerate}

\end{document}